% ----------------------------------------------------------
\chapter{Resultados e Discussões}\label{cap:resultados}
% ----------------------------------------------------------

\section{Descrição da precipitação}

a seca

\section{Dinâmica em grande escala}

Para ET1 

\subsection{Análise do Vento Zonal em 200 hPa e 850 hPa}

Segundo  

\section{Análise Modal}

\subsection{Análise Espacial de Ondas}

A propagação das ondas atmosféricas, durante a seca extrema no Sudeste do Brasil, como as ondas de Kelvin (KW

\subsection{Ondas de Gravidade Inerciais para Oeste (WIG)}

A Figu

\subsection{Ondas de Kelvin (KW)}

A Figuritação, contribuindo para um padrão de seca no Sudeste do Brasil. Na região do Pacífico Ocidental, também se observa uma prevalência da fase positiva, o que favorece a convecção nesta área.
